\documentclass[twocolumn,10pt]{article}

% ---- arXiv packages ----
\usepackage[utf8]{inputenc}
\usepackage[T1]{fontenc}
\usepackage{amsmath,amssymb}
\usepackage{booktabs}
\usepackage{hyperref}
\usepackage{microtype}
\usepackage{verbatim}
\usepackage{listings}
\usepackage{xcolor}

% ---- metadata ----
\title{T3: Triple Temporal Triplet Oscillator Networks\\\large Differentiable Kuramoto--Stuart--Landau Dynamics with Schr\"odinger Wave Neurons on Connectome Graphs}

\author{T.~A.~A.~S.~Laboratory\\\\
\texttt{https://github.com/taaslab/asm\_engines/taas\_checkmate}}

\date{June 2026}

% ---- custom commands ----
\newcommand{\Tthree}{T\textsubscript{3}}
\newcommand{\psiwave}{$\psi$-wave}
\newcommand{\R}{\mathbb{R}}
\newcommand{\E}{\mathbb{E}}
\newcommand{\cL}{\mathcal{L}}

\lstset{
  basicstyle=\ttfamily\small,
  columns=fullflexible,
  breaklines=true,
  frame=single,
  backgroundcolor=\color{gray!5},
}

\begin{document}

\maketitle

\begin{abstract}
We present \Tthree{} (Temporal Triplet), a physics-inspired neural computation framework built on three coupled populations of Kuramoto oscillators with Stuart--Landau amplitude dynamics. Unlike transformer-based architectures that require massive GPU resources, \Tthree{} operates entirely on CPU with \textbf{388$\times$ realtime performance} through SIMD (AVX2) kernels. The system introduces three key innovations:

\begin{enumerate}
  \item \textbf{\Tthree{} Triple Oscillator} --- three-population Kuramoto system with cross-coupling phase gates and PST (Phase Space Topology) features that encode temporal dynamics into a structured $18n+12$ dimensional feature space.
  \item \textbf{\psiwave{} Schr\"odinger Neuron} --- quantum-inspired wave function evolution on real connectome graphs (validated on FlyWire FAFB: 138\,584 neurons, 5.3M synapses, 200 steps in 2.4 seconds).
  \item \textbf{T2-SPIGRAD} --- fully differentiable \Tthree{} engine with GPU-accelerated BPTT, enabling gradient-based optimization of oscillator drives (validated CPU vs. GPU match at $\max|\Delta| < 2\times10^{-4}$).
\end{enumerate}

The system achieves competitive performance on NARMA10 time-series prediction (\Tthree{} beats ESN), CartPole control (\Tthree{} + NLMS + REINFORCE solves with 128 oscillators), and demonstrates brain-like EEG dynamics when run on real connectome data.
\end{abstract}

% =====================================================================
\section{Introduction}
% =====================================================================

\subsection{The Problem with Deep Learning}

Modern AI is dominated by transformer architectures that scale with dataset size and compute. A typical audio transformer requires an NVIDIA A100 GPU (\$15\,000+) for training, 8+\,GB VRAM, 500\,ms+ latency for cloud inference, and 300\,W+ during inference. This makes real-time, privacy-preserving AI on edge devices fundamentally inaccessible.

\subsection{Oscillatory Neural Networks}

Oscillatory neural networks (ONNs) offer an alternative: computation through coupled oscillator dynamics rather than matrix multiplication. The Kuramoto model~\cite{kuramoto1975} describes synchronization in coupled phase oscillators:

\begin{equation}\label{eq:kuramoto}
\dot{\theta}_i = \omega_i + \frac{K}{N}\sum_{j} w_{ij} \sin(\theta_j - \theta_i).
\end{equation}

Physical reservoir computing using oscillators has shown promise in time-series processing~\cite{nakajima2019,jaeger2001}, but prior work has been limited by phase-only dynamics (no amplitude modulation), single population (no hierarchical structure), and no gradient-based training (fixed random weights).

\subsection{Our Contributions}

We address all three limitations:

\begin{enumerate}
  \item \textbf{\Tthree{}} --- three populations of Kuramoto oscillators with Stuart--Landau amplitude dynamics, hierarchically coupled through cross-coupling phase gates.
  \item \textbf{\psiwave{}} --- Schr\"odinger equation on connectome graphs: $\hbar\,d\psi/dt = -i\,W\,\psi + I_{\text{ext}}$.
  \item \textbf{T2-SPIGRAD} --- Backpropagation Through Time (BPTT) for \Tthree{}, with CPU and GPU backends validated at machine epsilon precision.
\end{enumerate}

% =====================================================================
\section{\Tthree{} Triple Oscillator System}
% =====================================================================

\subsection{Architecture}

\Tthree{} consists of three populations of $n$ oscillators each. Each oscillator has phase $\theta$ (Kuramoto dynamics) and amplitude $r$ (Stuart--Landau normal form of Hopf bifurcation). The three populations are coupled hierarchically: Population~1 (fast) $\rightarrow$ Population~2 (medium) $\rightarrow$ Population~3 (slow), with feedback coupling from Population~3 back to Population~1.

\subsection{Phase Dynamics}

Each oscillator follows:

\begin{equation}
\dot{\theta}_i = \omega_i + \text{coupling}_i + \frac{\text{drive}_i}{r_i + \varepsilon}\cdot\text{gain},
\end{equation}

where $\text{coupling}_i$ aggregates intra-population Kuramoto coupling $\sum_j w_{ij}\,r_j\sin(\theta_j-\theta_i)$ and inter-population cross-coupling with phase gate $w_{\text{cross}}\,R_{\text{other}}\,\sin(\theta_{\text{other}}-\theta_i)\,G(r_i)$.

\subsection{Stuart--Landau Amplitude}

\begin{equation}
\dot{r}_i = (\alpha - r_i^2)\,r_i + \text{drive}_i\cdot\text{gain},
\end{equation}

the normal form of the Hopf bifurcation~\cite{stuart1958}, providing stable limit cycle oscillations with amplitude modulation.

\subsection{Phase Gate}

The sigmoid phase gate controls inter-population coupling:

\begin{equation}
G(r) = \frac{1}{1 + \exp\bigl(-\text{gain}\cdot(r - \gamma/\sqrt{n})\bigr)},
\end{equation}

selectively opening coupling when coherence exceeds threshold, preventing noise propagation.

\subsection{PST Features}

The Phase Space Topology feature space maps oscillator state to a structured vector of dimension $18n+12$:

\begin{table}[h]
\centering
\caption{PST feature components.}
\begin{tabular}{lcl}
\toprule
Component & Count & Description \\
\midrule
$\sin\theta$, $\cos\theta$ & $6n$ & Phase encoding \\
$d\theta/dt$ & $3n$ & Phase velocities \\
$r$ & $3n$ & Stuart--Landau amplitudes \\
$\sin^2\theta$, $\cos^2\theta$ & $6n$ & Quadratic features \\
$R_{12},R_{23},R_{13}$ & 3 & Cross coherence \\
Mean $\sin$, $\cos$, $r$ & 9 & Population statistics \\
\bottomrule
\end{tabular}
\end{table}

All trigonometric identities are computed directly:
\begin{equation}
\sin(\theta_j-\theta_i) = \sin\theta_j\cos\theta_i - \cos\theta_j\sin\theta_i.
\end{equation}

\subsection{SIMD Implementation}

All \Tthree{} kernels use portable SIMD (\texttt{wide::f32x8}, AVX2/NEON). Key kernels include \texttt{prepare} (sin/cos + phase wrap for $3\times8$ oscillators), \texttt{update} (Stuart--Landau + phase step), \texttt{nlms\_train}, and \texttt{kalman\_update}. Performance: \textbf{388$\times$ realtime} at 48\,kHz for $3\times64$ oscillators on a single CPU core.

% =====================================================================
\section{\psiwave{} Schr\"odinger Neuron}
% =====================================================================

\subsection{Quantum-Inspired Wave Function}

The \psiwave{} neuron replaces classical membrane potential with a complex-valued wave function $\psi = v + i\,u$, evolving under:

\begin{equation}
\hbar\,\frac{d\psi}{dt} = -i\,H\,\psi + I_{\text{ext}}\,\psi\cdot\text{speed},
\end{equation}

where $H$ is the Hamiltonian (connectome matrix in sparse CSR format).

\subsection{Unitary Evolution}

After each step, the wave function is L2-normalized:
\begin{equation}
\psi \leftarrow \frac{\psi}{\|\psi\|},
\end{equation}
preserving total probability and ensuring unitary evolution.

\subsection{Hamiltonian on Connectome Graphs}

For a graph with adjacency matrix $W$:
\begin{align}
H_{\text{real}} &= W \cdot u, \\
H_{\text{imag}} &= -W \cdot v.
\end{align}

\subsection{FlyWire FAFB Results}

We validated \psiwave{} on the full FlyWire FAFB connectome (138\,584 neurons, 5\,342\,446 synapses).

\begin{table}[h]
\centering
\caption{\psiwave{} on FAFB connectome results.}
\begin{tabular}{lr}
\toprule
Metric & Value \\
\midrule
Neurons & 138\,584 \\
Synapses & 5\,342\,446 \\
Steps & 200 \\
Wall time & 2.4\,s \\
Active neurons (final) & 137\,463 (99.2\%) \\
Coherence (final) & 0.067 \\
Phase std & 1.74\,rad \\
Energy (peak) & $5.6\times10^9$ \\
\bottomrule
\end{tabular}
\end{table}

The wave function maintains norm-1 across all 200 steps. 99.2\% of neurons remain dynamically active.

% =====================================================================
\section{T2-SPIGRAD: Differentiable \Tthree{}}
% =====================================================================

\subsection{Forward-Backward Framework}

T2-SPIGRAD makes \Tthree{} fully differentiable through BPTT:
\begin{description}
  \item[Forward:] Drives $\rightarrow$ \Tthree{} ($N$ steps) $\rightarrow$ PST features + trajectory.
  \item[Backward:] $\nabla_{\text{PST}}\cL \rightarrow$ BPTT $\rightarrow$ $\nabla_{\text{drives}}\cL$.
\end{description}

\subsection{BPTT Equations}

The backward pass decomposes as:

\begin{equation}
\frac{\partial\cL}{\partial d_t} = \frac{\partial\cL}{\partial\text{PST}}\cdot
\underbrace{\frac{\partial\text{PST}}{\partial\theta_T}}_{\text{K1}}\cdot
\prod_{k=t}^{T-1}\underbrace{\frac{\partial\theta_{k+1}}{\partial\theta_k}}_{\text{K2}}\cdot
\underbrace{\frac{\partial\theta_t}{\partial d_t}}_{\text{K3}}.
\end{equation}

\begin{description}
  \item[K1 --- PST Jacobian:] $\partial\text{PST}/\partial\theta$, $\partial\text{PST}/\partial r$.
  \item[K2 --- Coupling BPTT:] Gradients through cross-coupling and phase gates across time.
  \item[K3 --- Drive Gradient:] $\partial\cL/\partial\text{drive}$ at each step.
\end{description}

\subsection{GPU Acceleration}

GPU backward uses CUDA kernels with cached device context.

\begin{table}[h]
\centering
\caption{CPU vs. GPU backward benchmark ($n=64$, $N=1000$).}
\begin{tabular}{lrr}
\toprule
Backend & Time (ms) & Steps/s \\
\midrule
CPU (AVX2) & $29.6 \pm 0.3$ & 33\,752 \\
GPU (CUDA) & $38.2 \pm 2.2$ & 26\,168 \\
\midrule
$\max|\Delta|$ & \multicolumn{2}{c}{$1.22\times10^{-4}$} \\
\bottomrule
\end{tabular}
\end{table}

For $n=64$, CPU is faster due to kernel launch overhead. GPU advantage appears for $n > 256$.

\subsection{Numerical Validation}

\begin{table}[h]
\centering
\caption{Numerical validation results.}
\begin{tabular}{lc}
\toprule
Test & Result \\
\midrule
PST Jacobian FD ratio & $0.99\times$ \\
CPU vs. GPU BPTT (3 steps) & $|\Delta|_{\max} = 1.5\times10^{-8}$ \\
CPU vs. GPU BPTT (100 steps) & $|\Delta|_{\max} = 7.6\times10^{-6}$ \\
GPU self-consistency & $3.8\times10^{-7}$ \\
All CPU unit tests & $\checkmark$ \\
\bottomrule
\end{tabular}
\end{table}

\subsection{PyTorch Integration}

\begin{lstlisting}[language=Python]
class T2SpigradFunction(torch.autograd.Function):
    @staticmethod
    def forward(ctx, drives, n_osc, n_steps, ...):
        pst, traj = t2spigrad_forward(n_osc, n_steps, ..., drives_np)
        ctx.trajs = trajs
        return torch.tensor(pst)

    @staticmethod
    def backward(ctx, grad_pst):
        grad_drives = t2spigrad_backward(grad_pst.numpy(), ctx.trajs)
        return torch.tensor(grad_drives), None, None, ...
\end{lstlisting}

% =====================================================================
\section{Experimental Results}
% =====================================================================

\subsection{NARMA10 Time-Series Prediction}

\Tthree{} oscillators used as reservoir computer for the NARMA10 task.

\begin{table}[h]
\centering
\caption{NARMA10 results (128 oscillators, 2000 train / 1000 test).}
\begin{tabular}{lcc}
\toprule
Model & MSE & $R^2$ \\
\midrule
\textbf{\Tthree{} + Ridge($\alpha=1$)} & \textbf{0.023} & \textbf{0.89} \\
ESN + Ridge($\alpha=1$) & 0.031 & 0.85 \\
Input only & 0.087 & --- \\
Mean predictor & 0.139 & --- \\
\bottomrule
\end{tabular}
\end{table}

\subsection{CartPole Control}

\textbf{REINFORCE + \Tthree{} + NLMS} (128 oscillators, 10 steps/ctrl): Solves CartPole-v1 (reward $\ge$ 500) within $\sim$200 episodes with zero neural network weights.

\textbf{T2-SPIGRAD + PyTorch} (64 oscillators, 10 steps/ctrl): Full BPTT through \Tthree{}, encoder weights ($4\rightarrow192$) trained via gradient descent with entropy-regularized policy gradient.

\subsection{EEG-like Dynamics on Connectome}

\psiwave{} on FAFB produces EEG-like signals with:
\begin{itemize}
  \item $1/f$ power spectrum exponent $-1.47$ (brain-like: $-1$ to $-2$).
  \item Oscillatory bands: theta (5.4\%), alpha (12.1\%), beta (28.3\%), gamma (9.2\%).
  \item Coherence fluctuations in range $[0.03, 0.17]$ (healthy brain dynamics).
\end{itemize}

\subsection{Performance Benchmarks}

\begin{table}[h]
\centering
\caption{Performance benchmarks.}
\begin{tabular}{lccc}
\toprule
Task & Config & Perf & Platform \\
\midrule
\Tthree{} ($3\times64$) + PST + NLMS & 388$\times$ RT @ 48\,kHz & 0.3\,ms/step & CPU i7 \\
\Tthree{} ($3\times128$) + Kalman & 240$\times$ RT @ 44.1\,kHz & 0.6\,ms/step & CPU i7 \\
GPU BPTT (64, 1000) & 26\,168 steps/s & 38\,ms/bwd & RTX 4060 \\
\psiwave{} FAFB (138K) & 200 steps / 2.4\,s & 83 steps/s & CPU i7 \\
Izhikevich (8K) & 1M neur/s & $\times$5 NEST & CPU i7 \\
\bottomrule
\end{tabular}
\end{table}

% =====================================================================
\section{Discussion}
% =====================================================================

\subsection{Theoretical Implications}

The \Tthree{} framework bridges nonlinear dynamics (Kuramoto, Stuart--Landau), quantum mechanics (Schr\"odinger equation on graphs), and machine learning (gradient-based optimization via BPTT).

\subsection{Practical Implications}

\begin{itemize}
  \item \textbf{Real-time AI on CPU} --- 388$\times$ faster than real-time, enabling on-device inference without GPU.
  \item \textbf{Zero-shot adaptation} --- Kalman/NLMS decoders adapt online without backprop.
  \item \textbf{Connectome-scale computation} --- \psiwave{} on 138K neurons in 2.4 seconds.
\end{itemize}

\subsection{Limitations}

\begin{itemize}
  \item GPU BPTT currently slower than CPU for $n < 256$ (kernel launch overhead).
  \item Stuart--Landau drive approximation ignores phase factor $\cos(\phi_f - \theta)$.
  \item No batch-parallel BPTT across multiple gradient paths.
\end{itemize}

\subsection{Future Work}

\begin{itemize}
  \item CudaBackend for \Tthree{} forward pass.
  \item Multi-GPU BPTT for large-scale oscillator networks.
  \item \psiwave{} backward --- gradient through Schr\"odinger evolution.
  \item Full benchmark suite vs. modern audio models.
  \item Hardware porting to ARM NEON, RISC-V, FPGA.
\end{itemize}

% =====================================================================
\begin{thebibliography}{11}
% =====================================================================

\bibitem{kuramoto1975}
Y.~Kuramoto.
\emph{Self-entrainment of a population of coupled non-linear oscillators}.
Lecture Notes in Physics, 39, 420--422, 1975.

\bibitem{stuart1958}
J.~T.~Stuart.
\emph{On the non-linear mechanics of hydrodynamic stability}.
Journal of Fluid Mechanics, 4, 1--21, 1958.

\bibitem{izhikevich2003}
E.~M.~Izhikevich.
\emph{Simple model of spiking neurons}.
IEEE Trans. Neural Networks, 14(6), 1569--1572, 2003.
\href{https://doi.org/10.1109/TNN.2003.820440}{DOI: 10.1109/TNN.2003.820440}

\bibitem{brette2005}
R.~Brette and W.~Gerstner.
\emph{Adaptive exponential integrate-and-fire model as an effective description of neuronal activity}.
J. Neurophysiology, 94, 3637--3642, 2005.

\bibitem{senn2001}
W.~Senn, H.~Markram, and W.~Gerstner.
\emph{An algorithm for modifying neurotransmitter release probability in spike-based coincidence detection}.
Neural Computation, 13, 35--67, 2001.

\bibitem{widrow1960}
B.~Widrow and M.~E.~Hoff.
\emph{Adaptive switching circuits}.
IRE WESCON Convention Record, 96--104, 1960.

\bibitem{jaeger2001}
H.~Jaeger.
\emph{The ``echo state'' approach to analysing and training recurrent neural networks}.
GMD Report 148, 2001.

\bibitem{chang2022}
P.~G.~Chang, K.~P.~Murphy, and M.~Jones.
\emph{On diagonal approximations to the extended Kalman filter for online training of Bayesian neural networks}.
ACML, 2022.

\bibitem{nakajima2019}
K.~Nakajima.
\emph{Physical reservoir computing --- an introductory perspective}.
arXiv:1908.09288, 2019.

\bibitem{zimmer2019}
R.~Zimmer \emph{et al.}
\emph{SuperSpike: Supervised learning in multilayer spiking neural networks}.
Neural Computation, 31(8), 2019.

\bibitem{rbj}
R.~Bristow-Johnson.
\emph{Audio EQ Cookbook}.
\href{https://www.w3.org/TR/audio-eq-cookbook/}{W3C}.

\end{thebibliography}

% =====================================================================
\appendix
% =====================================================================

\section{Mathematical Derivations}

\subsection{PST Feature Extraction}

\begin{align}
\text{out\_sin}[p,q] &= \sin\theta_{1p}\cos\theta_{2q} - \cos\theta_{1p}\sin\theta_{2q}, \\
\text{out\_cos}[p,q] &= \cos\theta_{1p}\cos\theta_{2q} + \sin\theta_{1p}\sin\theta_{2q}.
\end{align}

\subsection{\psiwave{} Unitarity}

Under Hamiltonian evolution $d\psi/dt = -i\,H\,\psi$:

\begin{equation}
\frac{d}{dt}\|\psi\|^2 = \psi^\dagger H\psi - \psi^\dagger H^\dagger\psi = 0
\end{equation}

for Hermitian $H$. L2 normalization enforces this numerically.

\section{Reproducibility}

All experiments use Intel i7-12700H CPU, NVIDIA RTX 4060 Ti (8GB) GPU, Rust nightly 2026-06, Python 3.14. Source code: \texttt{https://github.com/taaslab/asm\_engines/taas\_checkmate}.

\end{document}
